\section{Interchanging direct products with tensor}
\label{section-products-tensor}

\noindent
Let $M$ be an $R$-module and let $(Q_{\alpha})_{\alpha \in A}$ be a family of
$R$-modules.  Then there is a canonical map $M \otimes_R \left( \prod_{\alpha
\in A} Q_{\alpha} \right) \to \prod_{\alpha \in A} ( M \otimes_R
Q_{\alpha})$ given on pure tensors by $x \otimes (q_{\alpha}) \mapsto (x
\otimes q_{\alpha})$.  This map is not necessarily injective or surjective, as
the following example shows.

\begin{example}
\label{example-Q-not-ML}
Take $R = \mathbf{Z}$, $M = \mathbf{Q}$, and consider the family $Q_n =
\mathbf{Z}/n$ for $n \geq 1$.  Then $\prod_n (M \otimes Q_n) = 0$.  However
there is an injection $\mathbf{Q} \to M \otimes (\prod_n Q_n)$
obtained by tensoring the injection $\mathbf{Z} \to \prod_n Q_n$ by
$M$, so $M \otimes (\prod_n Q_n)$ is nonzero.  Thus $M \otimes (\prod_n
Q_n) \to \prod_n (M \otimes Q_n)$ is not injective.

\medskip\noindent
On the other hand, take again $R = \mathbf{Z}$, $M = \mathbf{Q}$, and let $Q_n
= \mathbf{Z}$ for $n \geq 1$.  The image of $M \otimes (\prod_n Q_n)
\to \prod_n (M \otimes Q_n) = \prod_n M$ consists precisely of
sequences of the form $(a_n/m)_{n \geq 1}$ with $a_n \in \mathbf{Z}$ and $m$
some nonzero integer.  Hence the map is not surjective.
\end{example}

\noindent
We determine below the precise conditions needed on $M$ for the map $M
\otimes_R \left( \prod_{\alpha} Q_{\alpha} \right) \to \prod_{\alpha}
(M \otimes_R Q_{\alpha})$ to be surjective, bijective, or injective for all
choices of $(Q_{\alpha})_{\alpha \in A}$.  This is relevant because the modules
for which it is injective turn out to be exactly Mittag-Leffler modules
(Proposition \ref{proposition-ML-tensor}).  In what follows, if $M$ is an
$R$-module and $A$ a set, we write $M^A$ for the product $\prod_{\alpha \in A}
M$.

\begin{proposition}
\label{proposition-fg-tensor}
Let $M$ be an $R$-module.  The following are equivalent:
\begin{enumerate}
\item $M$ is finitely generated.
\item For every family $(Q_{\alpha})_{\alpha \in A}$ of $R$-modules, the
canonical map $M \otimes_R \left( \prod_{\alpha} Q_{\alpha} \right)
\to \prod_{\alpha} (M \otimes_R Q_{\alpha})$ is surjective.
\item For every $R$-module $Q$ and every set $A$, the canonical map $M
\otimes_R Q^{A} \to (M \otimes_R Q)^{A}$ is surjective.
\item For every set $A$, the canonical map $M \otimes_R R^{A} \to
M^{A}$ is surjective.
\end{enumerate}
\end{proposition}

\begin{proof}
First we prove (1) implies (2).  Choose a surjection $R^n \to M$ and
consider the commutative diagram
$$
\xymatrix{
R^n \otimes_R (\prod_{\alpha} Q_{\alpha})  \ar[r]^{\cong} \ar[d] &
\prod_{\alpha} (R^n \otimes_R Q_{\alpha}) \ar[d] \\
M \otimes_R (\prod_{\alpha} Q_{\alpha})  \ar[r] & \prod_{\alpha} ( M
\otimes_R Q_{\alpha}).
}
$$
The top arrow is an isomorphism and the vertical arrows are surjections.  We
conclude that the bottom arrow is a surjection.

\medskip\noindent
Obviously (2) implies (3) implies (4), so it remains to prove (4) implies (1).
In fact for (1) to hold it suffices that the element $d = (x)_{x \in M}$ of
$M^M$ is in the image of the map $f: M \otimes_R R^{M} \to M^M$.  In
this case $d = \sum_{i = 1}^{n} f(x_i \otimes a_i)$ for some $x_i \in M$ and
$a_i \in R^M$.  If for $x \in M$ we write $p_x: M^M \to M$ for the
projection onto the $x$-th factor, then
$$
x = p_x(d) = \sum\nolimits_{i = 1}^{n} p_x(f(x_i \otimes a_i)) =
\sum\nolimits_{i=1}^{n} p_x(a_i) x_i.
$$
Thus $x_1, \ldots, x_n$ generate $M$.
\end{proof}

\begin{proposition}
\label{proposition-fp-tensor}
Let $M$ be an $R$-module.  The following are equivalent:
\begin{enumerate}
\item $M$ is finitely presented.
\item For every family $(Q_{\alpha})_{\alpha \in A}$ of $R$-modules, the
canonical map $M \otimes_R \left( \prod_{\alpha} Q_{\alpha} \right)
\to \prod_{\alpha} (M \otimes_R Q_{\alpha})$ is bijective.
\item For every $R$-module $Q$ and every set $A$, the canonical map $M
\otimes_R Q^{A} \to (M \otimes_R Q)^{A}$ is bijective.
\item For every set $A$, the canonical map $M \otimes_R R^{A} \to
M^{A}$ is bijective.
\end{enumerate}
\end{proposition}

\begin{proof}
First we prove (1) implies (2).  Choose a presentation $R^m \to R^n
\to M$ and consider the commutative diagram
$$
\xymatrix{
R^m \otimes_R (\prod_{\alpha} Q_{\alpha}) \ar[r] \ar[d]^{\cong} & R^n
\otimes_R (\prod_{\alpha} Q_{\alpha}) \ar[r] \ar[d]^{\cong} & M \otimes_R
(\prod_{\alpha} Q_{\alpha}) \ar[r] \ar[d] & 0 \\
\prod_{\alpha} (R^m \otimes_R Q_{\alpha}) \ar[r] & \prod_{\alpha} (R^n
\otimes_R Q_{\alpha}) \ar[r] & \prod_{\alpha} (M \otimes_R Q_{\alpha})
\ar[r] & 0.
}
$$
The first two vertical arrows are isomorphisms and the rows are exact.  This
implies that the map
$M \otimes_R (\prod_{\alpha} Q_{\alpha}) \to
\prod_{\alpha} ( M \otimes_R Q_{\alpha})$
is surjective and, by a diagram chase, also injective.  Hence (2) holds.

\medskip\noindent
Obviously (2) implies (3) implies (4), so it remains to prove (4) implies (1).
From Proposition \ref{proposition-fg-tensor}, if (4) holds we already know that
$M$ is finitely generated.  So we can choose a surjection $F \to M$
where $F$ is free and finite.  Let $K$ be the kernel.  We must show $K$ is
finitely generated.  For any set $A$, we have a commutative diagram
$$
\xymatrix{
& K \otimes_R R^A \ar[r] \ar[d]_{f_3} & F \otimes_R R^A \ar[r]
\ar[d]_{f_2}^{\cong} & M \otimes_R R^A \ar[r] \ar[d]_{f_1}^{\cong} & 0 \\
0 \ar[r] & K^A \ar[r] & F^A \ar[r] & M^A \ar[r] & 0 .
}
$$
The map $f_1$ is an isomorphism by assumption, the map $f_2$ is an isomorphism
since $F$ is free and finite, and the rows are exact.  A diagram chase shows
that $f_3$ is surjective, hence by Proposition \ref{proposition-fg-tensor} we
get that $K$ is finitely generated.
\end{proof}

\noindent
We need the following lemma for the next proposition.

\begin{lemma}
\label{lemma-kernel-tensored-fp}
Let $M$ be an $R$-module, $P$ a finitely presented $R$-module, and $f: P
\to M$ a map.  Let $Q$ be an $R$-module and suppose $x \in \Ker(P
\otimes Q \to M \otimes Q)$.  Then there exists a finitely presented
$R$-module $P'$ and a map $f': P \to P'$ such that $f$ factors through
$f'$ and $x \in \Ker(P \otimes Q \to P' \otimes Q)$.
\end{lemma}

\begin{proof}
Write $M$ as a colimit $M = \colim_{i \in I} M_i$ of a directed system of
finitely presented modules $M_i$.  Since $P$ is finitely presented, the map $f:
P \to M$ factors through $M_j \to M$ for some $j \in I$.  Upon
tensoring by $Q$ we have a commutative diagram
$$
\xymatrix{
& M_j \otimes Q \ar[dr] & \\
P \otimes Q \ar[ur] \ar[rr] & & M \otimes Q .
}
$$
The image $y$ of $x$ in $M_j \otimes Q$ is in the kernel of $M_j \otimes Q
\to M \otimes Q$.  Since $M \otimes Q = \colim_{i \in I} (M_i
\otimes
Q)$, this means $y$ maps to $0$ in $M_{j'} \otimes Q$ for some $j' \geq j$.
Thus we may take $P' = M_{j'}$ and $f'$ to be the composite $P \to M_j
\to M_{j'}$.
\end{proof}

\begin{proposition}
\label{proposition-ML-tensor}
Let $M$ be an $R$-module.  The following are equivalent:
\begin{enumerate}
\item $M$ is Mittag-Leffler.
\item For every family $(Q_{\alpha})_{\alpha \in A}$ of $R$-modules, the
canonical map $M \otimes_R \left( \prod_{\alpha} Q_{\alpha} \right)
\to \prod_{\alpha} (M \otimes_R Q_{\alpha})$ is injective.
\end{enumerate}
\end{proposition}

\begin{proof}
First we prove (1) implies (2).  Suppose $M$ is Mittag-Leffler and let $x$ be
in the kernel of $M \otimes_R (\prod_{\alpha} Q_{\alpha}) \to
\prod_{\alpha} (M \otimes_R Q_{\alpha})$.  Write $M$ as a colimit $M =
\colim_{i \in I} M_i$ of a directed system of finitely presented modules
$M_i$.
 Then $M \otimes_R (\prod_{\alpha} Q_{\alpha})$ is the colimit of $M_i
\otimes_R (\prod_{\alpha} Q_{\alpha})$.  So $x$ is the image of an element
$x_i \in M_i \otimes_R (\prod_{\alpha} Q_{\alpha})$.  We must show that $x_i$
maps to $0$ in $M_j \otimes_R (\prod_{\alpha} Q_{\alpha})$ for some $j \geq
i$.  Since $M$ is Mittag-Leffler, we may choose $j \geq i$ such that $M_i
\to M_j$ and $M_i \to M$ dominate each other.  Then consider
the commutative diagram
$$
\xymatrix{
M \otimes_R (\prod_{\alpha} Q_{\alpha}) \ar[r] & \prod_{\alpha} (M
\otimes_R Q_{\alpha}) \\
M_i \otimes_R (\prod_{\alpha} Q_{\alpha}) \ar[r]^{\cong} \ar[d] \ar[u] &
\prod_{\alpha} (M_i \otimes_R Q_{\alpha}) \ar[d] \ar[u] \\
M_j \otimes_R (\prod_{\alpha} Q_{\alpha}) \ar[r]^{\cong} & \prod_{\alpha}
(M_j \otimes_R Q_{\alpha})
}
$$
whose bottom two horizontal maps are isomorphisms, according to Proposition
\ref{proposition-fp-tensor}.  Since $x_i$ maps to $0$ in $\prod_{\alpha} (M
\otimes_R Q_{\alpha})$, its image in $\prod_{\alpha} (M_i \otimes_R
Q_{\alpha})$ is in the kernel of the map $\prod_{\alpha} (M_i \otimes_R
Q_{\alpha}) \to \prod_{\alpha} (M \otimes_R Q_{\alpha})$.  But this
kernel equals the kernel of $\prod_{\alpha} (M_i \otimes_R Q_{\alpha})
\to \prod_{\alpha} (M_j \otimes_R Q_{\alpha})$ according to the
choice of $j$.  Thus $x_i$ maps to $0$ in $\prod_{\alpha} (M_j \otimes_R
Q_{\alpha})$ and hence to $0$ in $M_j \otimes_R (\prod_{\alpha} Q_{\alpha})$.

\medskip\noindent
Now suppose (2) holds. We prove $M$ satisfies formulation (1) of being
Mittag-Leffler from Proposition \ref{proposition-ML-characterization}.  Let $f:
P \to M$ be a map from a finitely presented module $P$ to $M$.  Choose
a set $B$ of representatives of the isomorphism classes of finitely presented
$R$-modules. Let $A$ be the set of pairs $(Q, x)$ where $Q \in B$ and $x \in
\Ker(P \otimes Q \to M \otimes Q)$.  For $\alpha = (Q, x) \in A$, we
write $Q_{\alpha}$ for $Q$ and $x_{\alpha}$ for $x$.  Consider the commutative
diagram
$$
\xymatrix{
M \otimes_R (\prod_{\alpha} Q_{\alpha}) \ar[r] &
\prod_{\alpha} (M \otimes_R Q_{\alpha}) \\
P \otimes_R (\prod_{\alpha} Q_{\alpha}) \ar[r]^{\cong} \ar[u] &
\prod_{\alpha} (P \otimes_R Q_{\alpha}) \ar[u] .
}
$$
The top arrow is an injection by assumption, and the bottom arrow is an
isomorphism by Proposition \ref{proposition-fp-tensor}.  Let $x \in P
\otimes_R (\prod_{\alpha} Q_{\alpha})$ be the element corresponding to
$(x_{\alpha}) \in \prod_{\alpha} (P \otimes_R Q_{\alpha})$ under this
isomorphism.  Then $x \in \Ker( P \otimes_R (\prod_{\alpha} Q_{\alpha})
\to M \otimes_R (\prod_{\alpha} Q_{\alpha}))$ since the top arrow in
the diagram is injective.  By Lemma \ref{lemma-kernel-tensored-fp}, we get a
finitely presented module $P'$ and a map $f': P \to P'$ such that $f: P
\to M$ factors through $f'$ and $x \in \Ker(P \otimes_R
(\prod_{\alpha} Q_{\alpha}) \to P' \otimes_R (\prod_{\alpha}
Q_{\alpha}))$.  We have a commutative diagram
$$
\xymatrix{
P' \otimes_R (\prod_{\alpha} Q_{\alpha}) \ar[r]^{\cong} &
\prod_{\alpha} (P' \otimes_R Q_{\alpha}) \\
P \otimes_R (\prod_{\alpha} Q_{\alpha}) \ar[r]^{\cong} \ar[u] &
\prod_{\alpha} (P \otimes_R Q_{\alpha}) \ar[u] .
}
$$
where both the top and bottom arrows are isomorphisms by Proposition
\ref{proposition-fp-tensor}.  Thus since $x$ is in the kernel of the left
vertical map, $(x_{\alpha})$ is in the kernel of the right vertical map.  This
means $x_{\alpha} \in \Ker(P \otimes_R Q_{\alpha} \to P' \otimes_R
Q_{\alpha})$ for every $\alpha \in A$.  By the definition of $A$ this means
$\Ker(P \otimes_R Q \to P' \otimes_R Q) \supset \Ker(P \otimes_R
Q \to M \otimes_R Q)$ for all finitely presented $Q$ and, since $f: P
\to M$ factors through $f': P \to P'$, actually equality holds.
 By Lemma \ref{lemma-domination-fp}, $f$ and $f'$ dominate each other.
\end{proof}

\begin{lemma}
\label{lemma-minimal-contains}
Let $M$ be a flat Mittag-Leffler module over $R$. Let $F$ be an $R$-module
and let $x \in F \otimes_R M$. Then there exists a smallest submodule
$F' \subset F$ such that $x \in F' \otimes_R M$.
Also, $F'$ is a finite $R$-module.
\end{lemma}

\begin{proof}
Since $M$ is flat we have $F' \otimes_R M \subset F \otimes_R M$
if $F' \subset F$ is a submodule, hence the statement makes sense.
Let $I = \{F' \subset F \mid x \in F' \otimes_R M\}$ and for
$i \in I$ denote $F_i \subset F$ the corresponding submodule.
Then $x$ maps to zero under the map
$$
F \otimes_R M \longrightarrow \prod (F/F_i \otimes_R M)
$$
whence by Proposition \ref{proposition-ML-tensor}
$x$ maps to zero under the map
$$
F \otimes_R M \longrightarrow \left(\prod F/F_i\right) \otimes_R M
$$
Since $M$ is flat the kernel of this arrow is
$(\bigcap F_i) \otimes_R M$ which proves that $F' = \bigcap F_i$.
To see that $F'$ is a finite module, suppose that
$x = \sum_{j = 1, \ldots, m} f_j \otimes m_j$ with $f_j \in F'$
and $m_j \in M$. Then $x \in F'' \otimes_R M$ where $F'' \subset F'$ is
the submodule generated by $f_1, \ldots, f_m$. Of course then
$F'' = F'$ and we conclude the final statement holds.
\end{proof}

\begin{lemma}
\label{lemma-pure-submodule-ML}
Let $0 \to M_1 \to M_2 \to M_3 \to 0$ be a
universally exact sequence of $R$-modules.  Then:
\begin{enumerate}
\item If $M_2$ is Mittag-Leffler, then $M_1$ is Mittag-Leffler.
\item If $M_1$ and $M_3$ are Mittag-Leffler, then $M_2$ is Mittag-Leffler.
\end{enumerate}
\end{lemma}

\begin{proof}
For any family $(Q_{\alpha})_{\alpha \in A}$ of $R$-modules we have a
commutative diagram
$$
\xymatrix{
0 \ar[r] & M_1 \otimes_R (\prod_{\alpha} Q_{\alpha}) \ar[r] \ar[d] & M_2
\otimes_R (\prod_{\alpha} Q_{\alpha}) \ar[r] \ar[d] & M_3 \otimes_R
(\prod_{\alpha} Q_{\alpha}) \ar[r] \ar[d] & 0 \\
0 \ar[r] & \prod_{\alpha}(M_1 \otimes Q_{\alpha}) \ar[r] & \prod_{\alpha}(M_2
\otimes Q_{\alpha}) \ar[r] & \prod_{\alpha}(M_3 \otimes Q_{\alpha})\ar[r] & 0
}
$$
with exact rows. Thus (1) and (2) follow from
Proposition \ref{proposition-ML-tensor}.
\end{proof}

\begin{lemma}
\label{lemma-quotient-module-ML}
Let $M_1 \to M_2 \to M_3 \to 0$ be an exact sequence of $R$-modules.
If $M_1$ is finitely generated and $M_2$ is Mittag-Leffler, then $M_3$
is Mittag-Leffler.
\end{lemma}

\begin{proof}
For any family $(Q_{\alpha})_{\alpha \in A}$ of $R$-modules,
since tensor product is right exact, we have a commutative diagram
$$
\xymatrix{
M_1 \otimes_R (\prod_{\alpha} Q_{\alpha}) \ar[r] \ar[d] & M_2
\otimes_R (\prod_{\alpha} Q_{\alpha}) \ar[r] \ar[d] & M_3 \otimes_R
(\prod_{\alpha} Q_{\alpha}) \ar[r] \ar[d] & 0 \\
\prod_{\alpha}(M_1 \otimes Q_{\alpha}) \ar[r] & \prod_{\alpha}(M_2
\otimes Q_{\alpha}) \ar[r] & \prod_{\alpha}(M_3 \otimes Q_{\alpha})\ar[r] & 0
}
$$
with exact rows. By Proposition \ref{proposition-fg-tensor}
the left vertical arrow is surjective. By
Proposition \ref{proposition-ML-tensor} the middle vertical arrow
is injective. A diagram chase shows the right vertical arrow
is injective. Hence $M_3$ is Mittag-Leffler by
Proposition \ref{proposition-ML-tensor}.
\end{proof}


\begin{lemma}
\label{lemma-colimit-universally-injective-ML}
If $M = \colim M_i$ is the colimit of a directed system of Mittag-Leffler
$R$-modules $M_i$ with universally injective transition maps, then $M$ is
Mittag-Leffler.
\end{lemma}

\begin{proof}
Let $(Q_{\alpha})_{\alpha \in A}$ be a family of $R$-modules. We have to
show that $M \otimes_R (\prod Q_\alpha) \to \prod M \otimes_R Q_\alpha$
is injective and we know that
$M_i \otimes_R (\prod Q_\alpha) \to \prod M_i \otimes_R Q_\alpha$
is injective for each $i$, see Proposition \ref{proposition-ML-tensor}.
Since $\otimes$ commutes with filtered colimits, it suffices to show that
$\prod M_i \otimes_R Q_\alpha \to \prod M \otimes_R Q_\alpha$
is injective. This is clear as each of the maps
$M_i \otimes_R Q_\alpha \to M \otimes_R Q_\alpha$ is injective
by our assumption that the transition maps are universally injective.
\end{proof}

\begin{lemma}
\label{lemma-direct-sum-ML}
If $M = \bigoplus_{i \in I} M_i$ is a direct sum of $R$-modules, then $M$ is
Mittag-Leffler if and only if each $M_i$ is Mittag-Leffler.
\end{lemma}

\begin{proof}
The ``only if'' direction follows from Lemma \ref{lemma-pure-submodule-ML} (1)
and the fact that a split short exact sequence is universally exact. The
converse follows from Lemma \ref{lemma-colimit-universally-injective-ML}
but we can also argue it directly as follows. First note that if $I$ is
finite then this follows from Lemma
\ref{lemma-pure-submodule-ML} (2).  For general $I$, if all $M_i$ are
Mittag-Leffler then we prove the same of $M$ by verifying condition (1) of
Proposition \ref{proposition-ML-characterization}.
Let $f: P \to M$ be a map from a finitely presented module $P$.
Then $f$ factors as
$P \xrightarrow{f'} \bigoplus_{i' \in I'} M_{i'} \hookrightarrow
\bigoplus_{i \in I} M_i$
for some finite subset $I'$ of $I$.  By the finite case
$\bigoplus_{i' \in I'} M_{i'}$ is Mittag-Leffler and hence there exists
a finitely presented  module $Q$ and a map $g: P \to Q$ such that $g$
and $f'$ dominate each other.  Then also $g$ and $f$ dominate each other.
\end{proof}

\begin{lemma}
\label{lemma-flat-ML-over-ML-ring}
Let $R \to S$ be a ring map. Let $M$ be an $S$-module.
If $S$ is Mittag-Leffler as an $R$-module, and $M$ is flat and Mittag-Leffler
as an $S$-module, then $M$ is Mittag-Leffler as an $R$-module.
\end{lemma}

\begin{proof}
We deduce this from the characterization of
Proposition \ref{proposition-ML-tensor}.
Namely, suppose that $Q_\alpha$ is a family of $R$-modules.
Consider the composition
$$
\xymatrix{
M \otimes_R \prod_\alpha Q_\alpha =
M \otimes_S S \otimes_R \prod_\alpha Q_\alpha \ar[d] \\
M \otimes_S \prod_\alpha (S \otimes_R Q_\alpha) \ar[d] \\
\prod_\alpha (M \otimes_S S \otimes_R Q_\alpha) =
\prod_\alpha (M \otimes_R Q_\alpha)
}
$$
The first arrow is injective as $M$ is flat over $S$ and $S$ is
Mittag-Leffler over $R$ and the second arrow is injective as
$M$ is Mittag-Leffler over $S$. Hence $M$ is Mittag-Leffler over $R$.
\end{proof}


